\subsubsection{The year is the same; the report is different}
\label{sec:comparative-report-editions}

An analyst asking for 2024 profit can legitimately encounter both 30 and
22.5. The first is what the company originally published; the second is the
corrected comparative in its later report. The useful-life error explains the
difference, but does not decide which number belongs in an analysis made at
a particular date. Each answer needs the identity of the report and the time
at which that report was available.

Keep the same synthetic asset, income and tax assumptions. To make the
information dates explicit, suppose the original 2024 report was authorized,
published and entered into the system on 31 March 2025. Management discovers
the error on 1 April and accepts its correction on 2 April. The 2025 annual
report is authorized on 31 March 2026, published on 2 April and entered into
the system on 3 April. These publication and ingestion dates are authored for
the example. They are not observations about a real issuer.

A latest-available query imposes two cutoffs. The public cutoff limits the
reports an outside analyst could obtain; the system cutoff limits the reports
actually recorded in the system. With a public cutoff of 1 April 2026, the
answer for 2024 still comes from the original report, even if the system is
queried much later. With a public cutoff of 2 April and a system cutoff of
2 April, the new report is public but has not yet entered the system, so the
same original report remains the answer. Once both cutoffs reach its respective
dates, the 2025 edition supplies corrected 2024 profit of 22.5. An explicit
request for the original edition can still retrieve 30 afterward. Discovery
of the error does not backdate publication of its correction.

Now reconstruct the positions shown in that later report. For this separate
historical source comparison, IAS 8.42--49 and IAS 1.40A--D in the retained
EU 2023/1803 text support the particular retrospective steps described here
\citep{eu20231803}. The corresponding IAS 12, IAS 16 and IAS 33 provisions
were checked for this example's tax, depreciation and private-entity EPS
assumptions. That comparison does not establish equivalence of entire
standards or the law applicable to a later actual filing.

The current reporting date is 31 December 2025 and the comparative date is
31 December 2024. Because the error also has a material effect on the opening
position of the comparative period, IAS 1.40A--B requires a third statement
at 1 January 2024. Materiality at that opening is a separate finding; the
mere existence of a restatement does not establish it. Starting with the
corrected 2023 close gives the following three positions, in currency units:

\begin{center}
\small
\begin{tabular}{lrrr}
\toprule
& 1 Jan. 2024 & 31 Dec. 2024 & 31 Dec. 2025\\
\midrule
Cash & 130 & 172.5 & 215\\
Net PPE & 100 & 80 & 60\\
Current tax payable & 7.5 & 7.5 & 7.5\\
Deferred tax liability & 0 & 0 & 0\\
Equity & 222.5 & 245 & 267.5\\
\bottomrule
\end{tabular}
\end{center}

Each column balances. From the earliest opening, annual profit of 22.5
raises equity first to 245 and then to 267.5. Cash rises by 42.5 each
year because receipts of 50 exceed payment of the preceding year's tax
by that amount. Depreciation of 20 explains why this cash increase exceeds
profit. Correcting earlier depreciation and deferred tax changes neither
those receipts nor the fiscal deductions.

The reconciliation to originally reported equity must preserve where each
part of the error arose. Original equity at the start of 2025 was 260.
The correction before 2024 reduces it by 7.5; the correction to 2024
profit reduces it by another 7.5. Only then is current profit added:
\[
 \underbrace{260}_{\text{original opening 2025}}
 -\underbrace{7.5}_{\text{error before 2024}}
 -\underbrace{7.5}_{\text{error during 2024}}
 +\underbrace{22.5}_{\text{2025 profit}}
 =267.5.
\]
The two corrections together are the opening-equity adjustment of 15
already explained above. Subtracting that 15 again would count the same
error twice. The report presents the corrected history; requesting it
does not create another journal or charge the prior error to 2025 profit.

The earliest opening also has a precise informational limit. Although it
equals the corrected 2023 closing position, it is not a full restated 2023
income statement in the 2025 edition. A query for the latest published
full-year 2023 column therefore still selects the 2024 edition, which
contains that earlier comparative. The return value must identify it as
original, even though later evidence establishes an error. A request for
a reconstructed 2023 income statement would ask a different question.
Keeping the period, edition and information cutoffs together prevents the
system from silently substituting the answer to that other question.
